\uuid{cE0h}
\titre{Bijection}
\auteur{bodin ; hanani ; mzari}
\organisation{exo7}
\licence{CC BY-NC-SA 4.0}
\datecreate{2020-01-20}
\niveau{L1}
\module{Analyse}
\chapitre{Dérivabilité des fonctions réelles}
\sousChapitre{Autre}
\difficulte{3}
\theme{Dérivées, bijection}
\exo7qcmid{lille-1.derivees.025}

\contenu{
\begin{question}
\qclassification{124.00}
\qtags{motcle=bijection}
Soit $\displaystyle f:\Rr\to \Rr$ telle que $f(x)=x^3+3x+1$. Quelles sont les bonnes réponses ?
\begin{answers}
    \good{$\forall x\in \Rr$, $f'(x)>0$}
    \good{$f$ est une bijection et $\displaystyle (f^{-1})'(1)=\frac{1}{3}$.}
    \bad{$f$ est une bijection et $\displaystyle (f^{-1})'(1)=\frac{1}{f'(1)}=\frac{1}{6}$.}
    \bad{$f$ est une bijection et $\displaystyle (f^{-1})'(x)=\frac{1}{f'(x)}$.}
\end{answers}
\begin{explanations}
On a $f'(x)=3x^2+3>0$ pour tout $x\in \Rr$. Ainsi $f$ est continue et est strictement croissante sur $\Rr$. Donc, d'après le théorème de la bijection, $f$ est une bijection et
$$\forall x\in \Rr,\; (f^{-1})'(x)=\frac{1}{f'\left(f^{-1}(x)\right)}.$$
En particulier, et puisque $f(0)=1$, $\displaystyle (f^{-1})'(1)=\frac{1}{f'(f^{-1}(1))}=\frac{1}{f'(0)}=\frac{1}{3}$.
\end{explanations}
\end{question}
}
